\subsection{Base Case: (\texttt{T.4}, \texttt{D.3}) -- Addition Case 1}
\[
\cinfer[T.4]{+(num[$n_1$]; num[$n_2$]): num}
  {\code{num[$n_1$]: num} & \code{num[$n_2$]: num}}
\]

\[
\cinfer[D.3]{+(num[$n_1$]; num[$n_2$]) $\stepto$ num[$n_1 + n_2$]}{ }
\]

Using rule \texttt{T.1}:

\[
\cinfer[T.1]{num[$n_1 + n_2$]: num}{ }
\qed
\]

\subsection{Inductive Case: (\texttt{T.4}, \texttt{D.1}) -- Addition Case 2}
\[
\cinfer[T.4]{+($e_1$; $e_2$): num}
  {\code{$e_1$: num} & \code{$e_2$: num}}
\qquad
\cinfer[D.1]{+($e_1$; $e_2$) $\stepto$ +($e'_1$; $e_2$)}{e_1 \stepto e'_1}
\]

We assume preservation holds for subexpressions.
Hence, by the inductive hypothesis,
\texttt{$e_1$: num} and $e_1 \stepto e'_1$ implies
\texttt{$e'_1$: num}.
Rule \texttt{T.4} gives us:

\[
\cinfer[T.4]{+($e'_1$; $e_2$): num}
  {\code{$e'_1$: num} & \code{$e_2$: num}}
\qed
\]

\subsection{Inductive Case: (\texttt{T.4}, \texttt{D.2}) -- Addition Case 3}
\[
\cinfer[T.4]{+(num[$n_1$]; $e_2$): num}
  {\code{num[$n_1$]: num} & \code{$e_2$: num}}
\qquad
\cinfer[D.2]{+(num[$n_1$]; $e_2$) $\stepto$ +(num[$n_1$]; $e'_2$)}{e_2 \stepto e'_2}
\]

Since \texttt{$e_2$: num} and $e_2 \stepto e'_2$,
by the inductive hypothesis, \texttt{$e'_2$: num}.

Rule \texttt{T.4} gives us:
\[
\cinfer[T.4]{+(num[$n_1$]; $e'_2$): num}
  {\cinfer[T.1]{num[$n_1$]: num}{ }
   &
   \code{$e'_2$: num}}
\qed
\]

\subsection{Concatenation Cases}
The proof of the concatenation cases are analogous
to the proofs of the addition cases by just replacing
the \code{num}'s with \code{str}'s and the
\code{+}'s with \code{\^{}}'s. $\qed$

\subsection{Substitution Lemma}
Before we can prove the next case for our preservation
proof, we need the Substitution Lemma.
Informally, this lemma states that 
we can substitute subexpressions that are of the same type in
an expression $e$ without changing the type of $e$.

\begin{lem}
\textbf{(Substitution)}
If $y' : \tau$ and $y : \tau \vdash e : \tau'$,
then $[y' / y]e : \tau'$.
\end{lem}

This lemma can be proved by a typical proof by
induction on the structure of $e$, so we skip
this proof for brevity.
Now we return back to the proof of preservation.

\subsection{Base Case: (\texttt{T.6}, \texttt{D.8}) -- Let Case 1}
\[
\cinfer[T.6]{let($x; e_1; e_2$): $\tau_2$}
  {e_1: \tau_1
   &
   x : \tau_1 \vdash e_2: \tau_2
  }
\]
\[
\cinfer[D.8]{let($x$; $e_1$; $e_2$) $\stepto$ $[e_1 / x]e_2$}{e_1 \; \texttt{value}}
\]

Since $e_1: \tau_1$ and $x : \tau_1 \vdash e_2: \tau_2$,
by substitution lemma, we have $[e_1 / x]e_2 : \tau_2$.
$\qed$

\subsection{Inductive Case: (\texttt{T.6}, \texttt{D.7}) -- Let Case 2}
\[
\cinfer[T.6]{let($x; e_1; e_2$): $\tau_2$}
  {e_1: \tau_1
   &
   x : \tau_1 \vdash e_2: \tau_2
  }
\]
\[
\cinfer[D.7]{let($x$; $e_1$; $e_2$) $\stepto$ let($x$; $e'_1$; $e_2$)}{e_1 \stepto e'_1}
\]

Since \texttt{$e_1$: $\tau_1$} and $e_1 \stepto e'_1$,
by the inductive hypothesis, \texttt{$e'_1$: $\tau_1$}.
Using rule T.6:

\[
\cinfer[T.6]{let($x; e'_1; e_2$): $\tau_2$}
  {e'_1: \tau_1
   &
   x : \tau_1 \vdash e_2: \tau_2
  }
\qed
\]
We have completed the proof of preservation for $\minilang$!

\subsection{Final Remarks of Preservation Proof}
We have proved preservation by showing that for
each possible combination of the typing and evaluation
judgments, preservation holds.
How does one know when such a combination is possible?
A combination is possible when there exists a unification
of the \emph{patterns} of the two judgments in question.
For example, notice in the preservation proof, there
was no case for typing rule \texttt{T.4} and evaluation
rule \texttt{D.4}.
The conclusion of rule \texttt{T.4} contains the expression
\code{+($e_1$; $e_2$)}.
We can think of the pattern of the expression \code{+($e_1$; $e_2$)}
as that same expression except the subexpressions or \emph{parameters}
$e_1$ and $e_2$ are variables that can be replaced with
any other expression.
In order for (\texttt{T.4}, \texttt{D.4}) to be a possible
combination case for the preservation proof, 
there must exists a unifier for the pattern of \code{+($e_1$; $e_2$)}
(conclusion expression of \texttt{T.4})
and the pattern of \code{\^{}($e_1$; $e_2$)}
(expression to the left of $\stepto$ in conclusion of \texttt{D.4}).
Because the first symbols in each of those two expressions contain
constant symbols (\code{+} and \code{\^{}}) that differ,
no such unification exists; so the (\texttt{T.4}, \texttt{D.4})
combination is not a possible case for the preservation proof.
These remarks hint at the fact that these type of proofs can be
automatically verified (e.g. by a proof assistant tool such as Twelf).

